\chapter{Introduction}
\chaptermark{Introduction}
\HRule\\[2cm]
\label{Chapter1}
\lhead{\emph{\textbf{Chapter 1:} Introduction}}

\begin{spacing}{1.6}

\section{Background and Motivation}

The proliferation of cloud computing has transformed how organisations analyse
and act on data. Machine learning models that once required expensive on-premise
infrastructure can now be accessed as cloud services, enabling smaller
organisations to leverage sophisticated analytics capabilities. Among these,
anomaly detection , the automated identification of observations that deviate
significantly from expected normal behaviour has emerged as a critical
capability. It underpins fraud detection in financial systems, intrusion
detection in network security, fault prediction in industrial operations, disease
surveillance in healthcare monitoring, and quality control in manufacturing
processes.

However, this shift to cloud-based anomaly detection introduces a severe and
often underappreciated privacy risk. Anomaly detection requires the cloud server
to access raw data in order to compute anomaly scores. In virtually every domain
where anomaly detection is useful, this data is highly sensitive:

\begin{itemize}
  \item In \textbf{healthcare}, continuous patient monitoring systems generate
  physiological signals --- heart rate, blood pressure, oxygen saturation,
  neurological activity that constitute protected health information. Exposure
  of this data to cloud providers violates patient confidentiality and, in many
  jurisdictions, applicable law.

  \item In \textbf{finance}, transaction records, account balances, and payment
  patterns reveal personal spending behaviour, business relationships, investment
  strategies, and operational details that competitors or adversaries could
  exploit.

  \item In \textbf{industrial operations}, process sensor readings encode
  operational parameters, equipment configurations, production rates, and
  process chemistry that represent valuable intellectual property and, in the
  context of critical infrastructure, potential national security assets.

  \item In \textbf{network security}, traffic patterns reveal internal
  infrastructure topology, communication behaviour, user activity, and the
  presence of sensitive systems that attackers could use to plan intrusions.
\end{itemize}

The fundamental tension is clear: \textit{effective anomaly detection requires
the server to see the data, but the data is too sensitive to share with the
server.} Conventional approaches to this tension are unsatisfying. Anonymisation
and aggregation reduce utility. Differential privacy adds noise that degrades
detection accuracy. Encrypting data during transmission (transport-layer
security) does not protect it at the server, which must decrypt before
processing.

\section{Homomorphic Encryption as a Solution}

Homomorphic Encryption (HE) offers a fundamentally different approach. Rather
than protecting data during transmission and exposing it at the server, HE
enables computation directly on ciphertext. Arithmetic operations like addition,
multiplication, and their compositions can be performed on encrypted values
such that the encrypted result, when decrypted by the data owner, equals the
result of the same operations on the original plaintexts.

This property has a profound consequence for privacy-preserving anomaly
detection: the detection server can compute anomaly scores on encrypted sensor
data without ever holding the decryption key or accessing any plaintext value.
The data owner (the client) encrypts their data using their public key, sends
the ciphertexts to the server, the server computes encrypted anomaly scores, and
returns them to the client, who decrypts using their secret key to obtain the
final anomaly classifications. At no point does the server observe or access any
individual data point.

The Cheon-Kim-Kim-Song (CKKS) scheme \cite{ckks}, introduced in 2017, is
particularly well suited to this application because it supports
\textit{approximate} arithmetic on real valued vectors. Unlike earlier schemes
that operated on integers, CKKS natively encodes floating point measurements in
ciphertext slots, making it directly applicable to the continuous sensor
readings and normalised feature vectors that characterise anomaly detection
datasets.

\section{The Central Technical Challenge: Division Under Encryption}

While HE enables addition and multiplication on ciphertexts, it does not
natively support division. This is a critical limitation for anomaly detection,
because the most effective unsupervised algorithms require division at
fundamental computational steps:

The Statistical Anomaly Detector (SAD) computes the per feature mean
$\mu_j = \frac{1}{N}\sum_{i=1}^N x_{ij}$ (requiring division by $N$) and the
variance normalised anomaly score $s(\mathbf{x}) = \sum_{j=1}^d
\frac{(x_j - \mu_j)^2}{\sigma_j^2}$ (requiring division by $\sigma_j^2$
for each feature).

PCA based anomaly detection similarly requires computing the feature mean
before centring the data for projection onto principal components.

Without division, these algorithms cannot be faithfully implemented under
homomorphic encryption. This thesis identifies this as the primary technical
obstacle and proposes a solution: the \textbf{multi-round interactive
protocol}, in which the encrypted denominator is sent to the key-holding data
owner, who decrypts it, computes the exact reciprocal in plaintext, and
re-encrypts the result and sends it back. This achieves zero approximation error at zero
additional HE depth cost, making it superior to both Newton-Raphson iterative
inversion (which diverges on realistic data distributions) and Chebyshev
polynomial approximation (which produces unacceptable errors at usable degree).

\section{Research Objectives}

This thesis pursues the following research objectives:

\begin{enumerate}
  \item \textbf{Empirical library selection:} Conduct a systematic benchmark
  of FHE libraries to identify the most suitable foundation for privacy-
  preserving anomaly detection, characterising the performance trade offs between
  arithmetic speed, multiplication depth, and hardware acceleration support.

  \item \textbf{Exact encrypted division:} Design, implement, and evaluate
  approaches to division under CKKS encryption, identifying the limitations of
  existing approximation methods and proposing an exact alternative.

  \item \textbf{HE-compatible algorithm design:} Implement Statistical Anomaly
  Detection and Principal Component Analysis entirely under CKKS encryption,
  designing the computation to minimise HE depth consumption and noise
  accumulation.

  \item \textbf{End-to-end pipeline:} Build a complete, dataset agnostic
  preprocessing and encryption pipeline that converts raw multivariate data to
  CKKS encrypted form suitable for server side anomaly scoring.

  \item \textbf{Network deployment:} Deploy the encrypted anomaly detection
  system as a client-server network application with real time monitoring,
  demonstrating practical usability beyond a research prototype.

  \item \textbf{Empirical evaluation:} Measure detection accuracy (AUC) of
  the HE-based system against plaintext baselines on real world benchmark
  datasets, quantifying the accuracy cost of privacy preserving encryption.
\end{enumerate}

\section{Research Gap}

A review of the existing literature on privacy-preserving machine learning
and anomaly detection reveals the following gaps that this thesis addresses:

\textbf{Gap 1 --- Transport vs computational privacy.} The overwhelming majority
of deployed ``privacy-preserving'' anomaly detection systems provide only
transport layer privacy: data is encrypted during network transmission but
decrypted at the server before processing. The server observes all plaintext
values. Genuinely privacy-preserving systems, where the server computes on
ciphertexts without ever decrypting individual observations, are rare.

\textbf{Gap 2 --- Supervised vs unsupervised methods.} The few existing systems
that do provide computational privacy are predominantly supervised: they require
labeled examples of anomalous behaviour for training. In practice, labeled
anomaly data is scarce, expensive to produce, and inherently incomplete, because
anomalies by definition include novel patterns not observed before. Unsupervised
methods, which learn only from normal behaviour, are far more practical but
are rarely implemented under homomorphic encryption.

\textbf{Gap 3 --- Division in HE-based ML.} No systematic treatment of the
encrypted division problem , which is fundamental to most anomaly detection
algorithms exists in the literature in the context of real valued CKKS based
anomaly detection. Existing works either avoid algorithms requiring division
or apply approximations without characterising their error on real data.

\textbf{Gap 4 --- End-to-end deployment.} Existing HE-based anomaly detection
research focuses on algorithm level demonstrations without network deployment,
practical batching strategies for large ciphertext transfers, or real time
monitoring infrastructure. This limits the transition from research prototype
to usable system.

\section{Contributions}

This thesis makes the following original contributions to the field of
privacy-preserving machine learning:

\begin{enumerate}
  \item \textbf{Multi-round interactive protocol for exact encrypted division.}
  A protocol that resolves the encrypted division problem exactly, with
  zero approximation error and zero additional HE depth cost. The protocol
  uses the key-holding data owner as a trusted computation agent for aggregate
  parameter inversion, which is shown to be privacy-preserving because only
  global model statistics (not individual data points) are revealed.

  \item \textbf{Systematic comparison of division approaches.}
  The first systematic comparison of Newton-Raphson iterative inversion,
  Chebyshev polynomial approximation, and the multi round protocol for encrypted
  division in the context of real valued anomaly detection, with quantitative
  characterisation of failure modes and error bounds on real data distributions.

  \item \textbf{HE-SAD: complete CKKS implementation of Statistical Anomaly
  Detection.} A 9-round homomorphic protocol implementing variance-normalised
  anomaly scoring entirely on CKKS ciphertexts. HE-SAD achieves AUC $= 0.9533$
  on the SWaT dataset, demonstrating that full CKKS privacy preserving
  implementation incurs no accuracy penalty compared to plaintext detection.

  \item \textbf{HE-PCA: complete CKKS implementation of PCA-based anomaly
  detection.} An 8-round homomorphic protocol implementing PCA reconstruction
  error scoring with automatic optimal component selection. The analysis of
  accuracy loss in HE-PCA versus plaintext PCA provides the first empirical
  characterisation of CKKS noise effects on matrix projection operations in the
  anomaly detection context.

  \item \textbf{Dataset-agnostic preprocessing and encryption pipeline.}
  A generalized 5-stage pipeline that converts any multivariate time-series
  dataset to CKKS-encrypted form, handling heterogeneous continuous and binary
  feature types, computing all normalization parameters from normal data only,
  and applying SIMD-optimized column-major and row-major encoding strategies
  that minimize noise accumulation during HE computation.

  \item \textbf{Network-deployed system with live monitoring.}
  A complete network deployment of the privacy-preserving detection system using
  FastAPI for the detection server, HTTP client with batched ciphertext transfer,
  asynchronous score computation, and a live SocketIO monitoring dashboard with
  browser based controls. The deployment demonstrates practical usability with
  only two terminal processes required.

  \item \textbf{FHE library benchmarking study.}
  A four experiment systematic benchmark of Microsoft SEAL, IBM HElib, and
  Pyfhel, covering arithmetic operation latency, rotation operation latency,
  maximum multiplication depth, and noise budget evolution, providing an
  empirical foundation for FHE library selection in ML workloads.
\end{enumerate}

\section{Scope and Assumptions}

This thesis operates under the following scope and assumptions:

\begin{itemize}
  \item The threat model assumes an \textit{honest-but-curious} server: the
  server follows the protocol correctly but attempts to infer information about
  the data from the ciphertexts and computation results it observes.

  \item The privacy guarantee covers individual data points. Aggregate model
  parameters ($N$, $\sigma^2_j$) are treated as public, analogous to the
  parameters of a deployed ML model.

  \item The evaluation focuses on the unsupervised anomaly detection setting,
  where only normal data is available for training.

  \item All experiments are conducted on a single machine (localhost) to
  characterize algorithmic performance independently of network conditions.

  \item CKKS is used without bootstrapping (leveled HE), relying on careful
  algorithm design to stay within the available multiplication budget and complexity.
\end{itemize}

\section{Thesis Organization}

The remainder of this thesis is organized as follows:

\textbf{Chapter 2} reviews the mathematical foundations of homomorphic
encryption, with particular attention to the CKKS scheme, its parameters, and
the trade-offs they entail. It covers the anomaly detection algorithms used,
related work on privacy-preserving ML and anomaly detection, and the benchmark
datasets used for evaluation.

\textbf{Chapter 3} presents the FHE library benchmarking study, detailing the
experimental methodology, results across four experimental categories, and the
empirical rationale for selecting Microsoft SEAL as the implementation
foundation.

\textbf{Chapter 4} describes the system architecture, the dataset agnostic
preprocessing and encryption pipeline, the CKKS encoding strategies, and the
multi-round interactive protocol for encrypted division, including a systematic
analysis of why alternative approaches fail.

\textbf{Chapter 5} presents the HE-SAD and HE-PCA algorithm implementations,
the network deployment architecture, implementation challenges and their
resolutions, plaintext baseline models, and experimental results comparing
privacy-preserving and plaintext anomaly detection.

\textbf{Chapter 6} discusses the key findings, limitations of the current
implementation, directions for future work, and concludes the thesis.

\end{spacing}
